\documentclass[border=4pt]{standalone}
\usepackage{amsmath,amssymb,mathtools,xcolor,varwidth}
\begin{document}
\begin{varwidth}{28em}
\large\bfseries
Here is the same device used in Lemma VII. We produce $AD$ to the remote point      $d$ and $AE$ to $e$, keeping $Ad : Ae = AD : AE$, and draw the ordinates $db$ and $ec$      parallel to the original $BD$ and $CE$. A similar curve $Abc$ is drawn out at      this fixed distance, and the tangent $Ag$ touches both curves at $A$, cutting      the four ordinates in $F, G, f, g$. The genius of the trick is scale invariance: while the      original figure near $A$ vanishes, this enlarged copy at $d, e$ stays a fixed finite size,      so we may take limits there in plain view. Because every length is proportional, the areas      $\triangle ABD$, $\triangle ACE$ stay in constant ratio to the blown-up areas $Abd$, $Ace$.
\end{varwidth}
\end{document}
